% [VERIFY] related work -- related_work.md RW01-RW12. Organized by method class (not paper-by-paper).
% Typed edges: refutes(highway), baseline(vgg/googlenet/prelu/bn-inception), imports(bn/msra/datasets/detectors),
% extends(residual repr, multigrid).
\section{Related Work}
\label{sec:related}

\noindent\textbf{Residual representations.}
The idea of learning with respect to residuals has a long history in vision and
numerical computation. In image recognition, encodings such as VLAD and the
Fisher Vector represent images by residual vectors with respect to a
dictionary, and encoding residuals has been shown to be more effective than
encoding original vectors for shallow retrieval and classification
pipelines~\cite{vlad2010,pq2011,fisher2007,vlfeat2008}. In low-level vision and graphics,
solving partial differential equations through multigrid methods or
hierarchical basis preconditioning reformulates a system in terms of residual
variables, and solvers that are aware of this residual nature converge far
faster than ones that are not~\cite{szeliski1990,szeliski2006,vatanen2013}. These results
suggest that a good reformulation or preconditioning can simplify optimization.
We extend this principle from hand-designed encodings and solvers to the
internal layers of deep networks, reformulating each block as a residual
function so the optimizer operates on residuals rather than unreferenced
mappings.

\noindent\textbf{Shortcut connections.}
Practices and theories that lead to shortcut connections have been studied for a
long time~\cite{bishop1995}. Early multi-layer perceptrons were trained with a linear layer
connecting the input to the output~\cite{ripley1996,venables1999}, and later
work centered layer responses, gradients, and propagated errors through
shortcut connections~\cite{raiko2012,schraudolph1998a,schraudolph1998b}.
Inception layers are also composed of a shortcut branch and a few deeper
branches~\cite{googlenet2015}. Concurrent with our work, ``highway networks''
present shortcut connections with gating
functions~\cite{highway2015,highwaydeep2015,lstm1997}. These gates are data-dependent and
have parameters, in contrast to our parameter-free identity shortcuts. When a
gated shortcut is ``closed'' (approaching zero), the layers in highway networks
represent non-residual functions; by contrast, our formulation always learns
residual functions, our identity shortcuts are never closed, and all
information is always passed through, with additional residual functions to be
learned. In addition, highway networks have not demonstrated accuracy gains with
extremely increased depth (e.g., over 100 layers).

\noindent\textbf{Deep network baselines and enabling techniques.}
VGG nets establish the design philosophy of stacking small $3\times3$
convolutions to great depth, and serve as the cleanest \plainnet{} baseline
against which residual learning is measured~\cite{vgg2015}. We adopt VGG's
design rules but operate at far greater depth and lower
complexity. Batch normalization is applied after each convolution and before
activation in both our \plainnet{}s and residual nets; it ensures forward
signals maintain non-zero variances and lets us isolate the
optimization-difficulty hypothesis for the \degrad{}~\cite{bn2015}. All
convolutional weights, in both \plainnet{} and residual variants, are
initialized with the MSRA scheme for rectifier
networks~\cite{msra2015}. Auxiliary-classifier and dropout-based regularizers
used by earlier deep models~\cite{dsn2014,dropout2012} are deliberately omitted
so that the depth ablation remains clean. We further compare our single-model
and ensemble ImageNet results against GoogLeNet, PReLU-net, and BN-inception as
strong contemporary reference points~\cite{googlenet2015,msra2015,bn2015}.

\noindent\textbf{Transfer to detection.}
To study whether deeper representations help beyond classification, we build on
the Faster R-CNN detection framework, which couples a region proposal network
with a Fast R-CNN detector~\cite{fasterrcnn2015,fastrcnn2015,rcnn2014}. Because
our 101-layer residual net lacks the hidden fully-connected layers that the VGG
backbone provides, we adopt the strategy of treating the convolutional feature
maps as shared full-image features and the final convolutional stage as the
per-region equivalent of the fully-connected
layers~\cite{noc2015,spp2014}. Holding the detection hyperparameters fixed and
swapping only the backbone isolates the contribution of the learned
representation. We evaluate on the standard PASCAL VOC and COCO detection
benchmarks~\cite{pascalvoc2010,coco2014,imagenet2015}.
